\documentclass[conference]{IEEEtran}
\IEEEoverridecommandlockouts

% ---------------------------------------------------------------
% Packages
% ---------------------------------------------------------------
\usepackage[utf8]{inputenc}
\usepackage[T1]{fontenc}
\usepackage[french]{babel}
\usepackage{cite}
\usepackage{amsmath,amssymb,amsfonts}
\usepackage{algorithmic}
\usepackage{graphicx}
\usepackage{textcomp}
\usepackage{xcolor}
\usepackage{booktabs}
\usepackage{hyperref}
\usepackage{multirow}

% ---------------------------------------------------------------
% Titre et auteurs
% ---------------------------------------------------------------
\title{Exploration d'architectures neuronales profondes pour l'évaluation du risque ordinal}

\author{
  \IEEEauthorblockN{Noé Cochet}
  \IEEEauthorblockA{\textit{Dép. d'informatique et de génie logiciel}\\
  Université Laval, Québec, Canada\\
  noe-pierre.cochet.1@ulaval.ca}
  \and
  \IEEEauthorblockN{Théo Fornier}
  \IEEEauthorblockA{\textit{Dép. d'informatique et de génie logiciel}\\
  Université Laval, Québec, Canada\\
  theo.fornier.1@ulaval.ca}
  \and
  \IEEEauthorblockN{Bilal Kefif}
  \IEEEauthorblockA{\textit{Dép. d'informatique et de génie logiciel}\\
  Université Laval, Québec, Canada\\
  bilal.kefif.1@ulaval.ca}
  \and
  \IEEEauthorblockN{Étienne Plante-Dubé}
  \IEEEauthorblockA{\textit{Dép. d'informatique et de génie logiciel}\\
  Université Laval, Québec, Canada\\
  etienne.plante-dube.1@ulaval.ca}
  \and
  \IEEEauthorblockN{Jean-Simon Ratté}
  \IEEEauthorblockA{\textit{Dép. d'informatique et de génie logiciel}\\
  Université Laval, Québec, Canada\\
  jean-simon.ratte.1@ulaval.ca}
  }

\begin{document}

\maketitle

% ---------------------------------------------------------------
% Résumé
% ---------------------------------------------------------------
\begin{abstract}
% TODO : Rédiger le résumé en dernier, une fois les résultats obtenus.
% Environ 150–250 mots. Doit couvrir : contexte, problème, approche, résultats principaux, conclusion.
Ce travail explore l'application de méthodes d'apprentissage profond récentes au problème de
classification ordinale du risque en assurance-vie, en utilisant le jeu de données
\textit{Prudential Life Insurance Assessment}. Ce dataset tabulaire de 128 variables présente
un défi représentatif d'un domaine traditionnellement dominé par les méthodes de
\textit{gradient boosting}. Nous comparons des modèles de référence basés sur XGBoost,
LightGBM et CatBoost à des architectures d'apprentissage profond récentes conçues pour les
données tabulaires, notamment TabPFN-2.5, TabM et TabDPT. Plusieurs variantes
expérimentales sont étudiées~: utilisation de données non étiquetées par pseudo-étiquetage,
régularisation (Dropout, Batch Normalization, Weight Decay), et une approche
\textit{teacher--student} pour l'augmentation tabulaire synthétique. Une étude d'ablation
permet d'isoler la contribution des composantes clés de l'architecture sélectionnée.
Les résultats montrent que \textbf{[RÉSUMÉ DES RÉSULTATS À COMPLÉTER]}.
\end{abstract}

\begin{IEEEkeywords}
apprentissage profond, données tabulaires, classification ordinale, assurance-vie, semi-supervisé, transformer, gradient boosting
\end{IEEEkeywords}

% ---------------------------------------------------------------
% 1. Introduction
% ---------------------------------------------------------------
\section{Introduction}

L'évaluation du risque lors de la souscription en assurance-vie repose traditionnellement sur des longs questionnaires médicaux afin de éter. Depuis quelques années, les assureurs ont introduit des algorithmes de machine learning sur les questionnaires médicaux afin de réduire les commandes d'exigence médicales tout en minimisant le risque additionnel de mortalité.
notation complexes intégrant de nombreuses variables médicales, comportementales et
démographiques. Ce processus, qui vise à attribuer une classe de risque ordinal à chaque
client, constitue un problème de classification ordinale à huit niveaux.

Sur ce type de données tabulaires structurées, les méthodes de \textit{gradient boosting}
telles que XGBoost~\cite{chen2016xgboost}, LightGBM~\cite{ke2017lightgbm} et CatBoost~\cite{prokhorenkova2018catboost} demeurent l'état de l'art. Pourtant, l'essor récent
d'architectures de réseaux de neurones profonds conçues spécifiquement pour les données
tabulaires, comme les Transformers tabulaires et les modèles de fondation, laisse
entrevoir un potentiel compétitif croissant.

Ce projet réalisé dans le cadre du cours GLO-7030 vise à évaluer dans quelle
mesure les avancées récentes en apprentissage profond peuvent rivaliser avec  les
modèles de boosting classiques sur un problème réel de risque de mortalité en assurance-vie. Notre objectif principal est de comparer les architectures modernes d'apprentissage profond sur des données tabulaires (TabPFN-2.5, TabM, TabDPT) combinée à une étude d'ablation sur les composantes d'entraînement, ainsi qu'une
exploration de stratégies semi-supervisées et d'augmentation de données synthétiques.

La suite de ce rapport est organisée comme suit~: la Section~\ref{sec:etat_art} présente
l'état de l'art; la Section~\ref{sec:approche} décrit l'approche proposée; la
Section~\ref{sec:methodo} détaille la méthodologie et les expérimentations; la
Section~\ref{sec:resultats} présente et discute les résultats; la Section~\ref{sec:conclusion}
conclut le rapport.

% ---------------------------------------------------------------
% 2. État de l'art
% ---------------------------------------------------------------
\section{État de l'art}
\label{sec:etat_art}

\subsection{Gradient Boosting pour les données tabulaires}

Les méthodes de \textit{gradient boosting} constituent depuis plusieurs années l'approche
dominante pour les tâches d'apprentissage supervisé sur données tabulaires. XGBoost~\cite{chen2016xgboost}
introduit un cadre d'optimisation scalable avec régularisation L1/L2 intégrée. LightGBM~\cite{ke2017lightgbm}
améliore l'efficacité computationnelle grâce à une construction d'arbres feuille par feuille
et un échantillonnage par gradient. CatBoost~\cite{prokhorenkova2018catboost} propose un
traitement natif des variables catégorielles via un encodage ordonné.

\subsection{Apprentissage profond sur des données tabulaires}

Plusieurs travaux récents ont tenté de combler l'écart de performance entre réseaux de
neurones et méthodes de boosting sur les données tabulaires~\cite{grinsztajn2022tree}.

\subsubsection{Modèles de fondation tabulaires}
TabPFN~\cite{hollmann2022tabpfn} est un modèle de fondation \textit{prior-fitted} entraîné
sur des millions de jeux de données synthétiques, permettant une inférence en contexte sans
réentraînement. La version 2.5 étendue améliore la scalabilité et la gestion de la haute
dimensionnalité~\cite{hollmann2025tabpfn2}.

\subsubsection{Approches d'ensemble et de mélange d'experts}
TabM~\cite{gorishniy2024tabm} combine des réseaux MLP avec un mécanisme de mélange
d'experts pondérés pour mieux capturer les interactions entre variables.

\subsubsection{Transformers semi-supervisés}
TabDPT~\cite{tabdpt2024} adapte l'architecture Transformer à l'apprentissage tabulaire semi-supervisé,
exploitant des données non étiquetées pour améliorer les représentations apprises.

\subsection{Classification ordinale}

La prédiction d'une variable cible ordonnée nécessite des ajustements par rapport à la
classification multiclasse standard. Parmi les approches courantes figurent la
\textit{ordinal regression}~\cite{frank2001simple}, la décomposition en classifieurs
binaires cumulatifs, et les fonctions de perte tenant compte de l'ordre (ex.~: \textit{Earth
Mover's Distance}, loss d'ordonnancement)~\cite{TODO}.

\subsection{Apprentissage semi-supervisé et augmentation tabulaire}

L'apprentissage semi-supervisé pour les données tabulaires peut prendre la forme de
pseudo-étiquetage~\cite{lee2013pseudo} ou de pré-entraînement par auto-encodage~\cite{TODO}.
L'approche \textit{teacher--student}~\cite{hinton2015distilling} a montré son efficacité pour
transférer la connaissance d'un modèle supervisé vers un modèle entraîné sur des données
synthétiques ou non étiquetées.

% ---------------------------------------------------------------
% 3. Approche proposée
% ---------------------------------------------------------------
\section{Approche proposée}
\label{sec:approche}

\subsection{Vue d'ensemble}

Notre approche suit un cadre comparatif structuré en deux volets~: (1) l'établissement de
lignes de base avec des modèles de \textit{gradient boosting}, et (2) l'exploration
d'architectures d'apprentissage profond récentes adaptées aux données tabulaires, avec
plusieurs variantes expérimentales.

\subsection{Description du jeu de données}

Le jeu de données \textit{Prudential Life Insurance Assessment}~\cite{kaggle_prudential}
est composé de \textbf{[N\_TRAIN]} exemples d'entraînement et \textbf{[N\_TEST]} exemples
de test, chacun décrit par 128 variables incluant des caractéristiques médicales, comportementales
et démographiques. La variable cible \texttt{Response} prend 8 valeurs ordinales ($1$ à $8$).
La distribution des classes est \textbf{[DÉCRIRE DÉSÉQUILIBRE]}.

Les variables peuvent être regroupées en~:
\begin{itemize}
  \item Variables continues (ex.~: \texttt{Ht}, \texttt{Wt}, \texttt{BMI})~;
  \item Variables catégorielles à haute cardinalité (ex.~: \texttt{Product\_Info\_2})~;
  \item Variables binaires médicales (ex.~: \texttt{Medical\_History\_*})~;
  \item Variables avec valeurs manquantes (ex.~: \texttt{Employment\_Info\_*}).
\end{itemize}

\subsection{Prétraitement}

\textbf{[DÉCRIRE LE PRÉTRAITEMENT]} : normalisation, encodage des variables catégorielles
(ex.~: ordinal encoding, target encoding, embeddings), imputation des valeurs manquantes,
gestion des variables fortement corrélées.

\subsection{Modèles de référence (\textit{Gradient Boosting})}

Nous entraînons trois modèles de référence avec optimisation des hyperparamètres par
validation croisée~:
\begin{itemize}
  \item \textbf{XGBoost}~\cite{chen2016xgboost}~: \textbf{[HYPERPARAMÈTRES CLÉS]}~;
  \item \textbf{LightGBM}~\cite{ke2017lightgbm}~: \textbf{[HYPERPARAMÈTRES CLÉS]}~;
  \item \textbf{CatBoost}~\cite{prokhorenkova2018catboost}~: \textbf{[HYPERPARAMÈTRES CLÉS]}.
\end{itemize}

\subsection{Architectures d'apprentissage profond}

\subsubsection{TabPFN-2.5}
\textbf{[DÉCRIRE L'UTILISATION DE TABPFN-2.5]} : configuration, adaptation au problème
ordinal, limites de scalabilité rencontrées.

\subsubsection{TabM}
TabM~\cite{gorishniy2024tabm} repose sur un mécanisme de \textit{BatchEnsemble} appliqué
à chaque couche linéaire d'un MLP standard. Pour chaque couche, $k$ adaptateurs légers
(scalaires d'entrée $r_i$, de sortie $s_i$ et de biais $b_i$, avec $i \in \{1,\ldots,k\}$)
transforment le flux d'information de façon indépendante, créant ainsi $k$ représentations
parallèles qui partagent néanmoins les mêmes poids principaux. Cette architecture permet
d'obtenir la diversité d'un ensemble de $k$ modèles tout en conservant un coût
paramétrique quasi-identique à un MLP seul.

Nous explorons quatre variantes issues de l'article~:
\begin{itemize}
  \item \textbf{plain} (MLP de base, sans ensemble)~;
  \item \textbf{tabm-packed} (ensemble de $k$ MLPs indépendants, sans partage de poids)~;
  \item \textbf{tabm-mini} (adapteur $R$ uniquement, un seul scalaire par membre)~;
  \item \textbf{tabm} (adapteurs $R$, $S$ et $B$ combinés, variante complète).
\end{itemize}

Les variantes notées TabM$^\dagger$ dans l'article ajoutent des embeddings
\textit{PiecewiseLinear}~\cite{gorishniy2022embeddings} en entrée, qui discrétisent chaque
variable continue en $B$ intervalles avant de la projeter dans un espace continu de
dimension $d_e$, améliorant significativement la modélisation des non-linéarités.

\subsubsection{TabDPT}
\textbf{[DÉCRIRE L'UTILISATION DE TABDPT]} : exploitation des données non étiquetées,
stratégie de pré-entraînement.

\subsection{Gestion de l'ordinalité}

Pour adapter les modèles de Deep Learning à la nature ordinale de la cible, nous utilisons
\textbf{[DÉCRIRE STRATÉGIE~: ex.~ordinal regression, coral loss, EMD loss, etc.]}.

% ---------------------------------------------------------------
% 4. Méthodologie et expérimentation
% ---------------------------------------------------------------
\section{Méthodologie et expérimentation}
\label{sec:methodo}

\subsection{Protocole d'évaluation}

Tous les modèles sont évalués selon les métriques suivantes~:
\begin{itemize}
  \item \textbf{Cohen's Kappa quadratique pondéré (QWK)}~: métrique principale du concours Kaggle~;
  \item \textbf{Accuracy}~: taux de classification exacte~;
  \item \textbf{MAE ordinal}~: erreur absolue moyenne sur les niveaux de risque.
\end{itemize}

Nous utilisons une validation croisée stratifiée à $k=5$ plis pour tous les modèles supervisés.
Les résultats sur le jeu de test Kaggle sont obtenus via soumission.

\subsection{Détails d'entraînement}

\textbf{[COMPLÉTER POUR CHAQUE MODÈLE]}~:
\begin{itemize}
  \item Matériel utilisé (GPU/CPU, mémoire)~;
  \item Optimiseur (ex.~: Adam, AdamW), taux d'apprentissage, scheduler~;
  \item Nombre d'époques, \textit{early stopping}~;
  \item Taille de lot (\textit{batch size})~;
  \item Initialisation des poids.
\end{itemize}

\subsection{Variantes expérimentales}

\subsubsection{Expérience 1 : Comparaison des modèles de base}
Comparaison directe de tous les modèles dans leur configuration standard sur les données
d'entraînement avec prétraitement minimal.

\subsubsection{Expérience 2 : Régularisation}
Évaluation de l'impact de différentes techniques de régularisation sur le modèle DL retenu~:
Dropout ($p \in \{0.1, 0.3, 0.5\}$), Batch Normalization, Weight Decay
($\lambda \in \{10^{-4}, 10^{-3}, 10^{-2}\}$).

\subsubsection{Expérience 3 : Utilisation des données non étiquetées}
Exploitation des données de test (sans étiquettes) via pseudo-étiquetage et/ou pré-entraînement
par auto-encodage pour améliorer les représentations apprises.

\subsubsection{Expérience 4 : Approche \textit{Teacher--Student} avec augmentation synthétique}
Un modèle enseignant (\textit{teacher}) entraîné sur les données réelles génère des
pseudo-étiquettes sur un grand volume de données synthétiques. Le modèle élève (\textit{student})
est pré-entraîné sur ces données augmentées, puis affiné (\textit{fine-tuned}) sur les données
originales.

\subsubsection{Expérience 5 : Étude d'ablation TabM}
Nous isolons la contribution de chaque composante de TabM sur le jeu de données Prudential.
Six variantes sont entraînées dans des conditions identiques~: même optimiseur AdamW,
scheduler CosineAnnealingLR, \textit{early stopping} à patience~$=10$, $k=32$ membres,
3~blocs de dimension~512 (2~blocs pour les variantes avec embeddings, conformément à
l'article). Un sweep bayésien via Weights \& Biases explore ensuite l'espace des
hyperparamètres de la meilleure variante (30~configurations, optimisation du QWK).
Enfin, un \textit{stacking} combinant TabM$^\dagger$, XGBoost et LightGBM est évalué
en validation croisée à 5~plis.

% ---------------------------------------------------------------
% 5. Résultats et Discussion
% ---------------------------------------------------------------
\section{Résultats et Discussion}
\label{sec:resultats}

\subsection{Comparaison globale des modèles}

\begin{table}[htbp]
\caption{Comparaison des performances des modèles sur la validation croisée ($k=5$)}
\label{tab:resultats_globaux}
\centering
\begin{tabular}{lccc}
\toprule
\textbf{Modèle} & \textbf{QWK} & \textbf{Accuracy} & \textbf{MAE} \\
\midrule
XGBoost          & -- & -- & -- \\
LightGBM         & -- & -- & -- \\
CatBoost         & -- & -- & -- \\
TabPFN-2.5       & -- & -- & -- \\
TabM$^\dagger$   & 0.630 & \textbf{0.6608} & -- \\
TabDPT           & -- & -- & -- \\
\bottomrule
\end{tabular}
\end{table}

TabM$^\dagger$ (avec embeddings PiecewiseLinear, optimisation bayésienne sur 30 runs W\&B)
atteint un QWK de \textbf{0,6608} sur l'ensemble de validation, soit le meilleur score
toutes architectures confondues. À titre de comparaison, le MLP de référence plafonne à
0,632. L'apprentissage profond spécialisé pour les données tabulaires (TabM) rivalise donc
avec les meilleures méthodes de \textit{gradient boosting} sur ce problème de régression ordinale.

\subsection{Impact de la régularisation}

\begin{table}[htbp]
\caption{Impact des techniques de régularisation sur le modèle retenu}
\label{tab:regularisation}
\centering
\begin{tabular}{lcc}
\toprule
\textbf{Configuration} & \textbf{QWK (val.)} & \textbf{QWK (test)} \\
\midrule
Sans régularisation & -- & -- \\
Dropout ($p=0.1$)   & -- & -- \\
Dropout ($p=0.3$)   & -- & -- \\
Batch Normalization & -- & -- \\
Weight Decay ($10^{-3}$) & -- & -- \\
\bottomrule
\end{tabular}
\end{table}

\subsection{Semi-supervisé et données non étiquetées}

\textbf{[DÉCRIRE ET ANALYSER LES RÉSULTATS]} : L'ajout des données non étiquetées a-t-il
amélioré les performances? Dans quelle mesure?

\subsection{Approche Teacher--Student}

\textbf{[DÉCRIRE ET ANALYSER LES RÉSULTATS]} : L'augmentation synthétique a-t-elle apporté
un gain? Comparer avec et sans pré-entraînement.

\subsection{Étude d'ablation TabM}

\subsubsection{Protocole}
Six variantes de TabM sont entraînées sur \textit{Prudential} ($N=59\,381$,
$d=130$~variables numériques après nettoyage) avec une division
80\,\%\,/\,20\,\% entraînement\,/\,validation stratifiée. L'optimiseur est
AdamW ($\eta=10^{-3}$, $\lambda=10^{-3}$), avec scheduler
\texttt{CosineAnnealingLR} et \textit{early stopping} à patience~$=10$.
La taille de lot est~$1024$; les variantes $\dagger$ utilisent $B=48$ intervalles
et $d_e=16$ (embeddings PiecewiseLinear). Les prédictions continues sont post-traitées
par optimisation d'offsets par classe, qui maximise le QWK sans ré-entraînement.

\subsubsection{Résultats}

\begin{table}[htbp]
\caption{Résultats de l'ablation TabM sur \textit{Prudential} (validation)}
\label{tab:ablation_tabm}
\centering
\begin{tabular}{lcccc}
\toprule
\textbf{Variante} & \textbf{arch} & \textbf{Embed.} &
\textbf{QWK\textsubscript{brut}} & \textbf{QWK\textsubscript{+off.}} \\
\midrule
MLP plain (baseline)    & plain      & --  & 0.594 & 0.632 \\
TabM-packed             & tabm-packed & --  & 0.593 & 0.637 \\
TabM-mini               & tabm-mini  & --  & 0.593 & 0.645 \\
TabM                    & tabm       & --  & 0.607 & 0.645 \\
TabM-mini$^\dagger$     & tabm-mini  & PL  & 0.608 & 0.656 \\
\textbf{TabM$^\dagger$} & \textbf{tabm} & \textbf{PL} &
\textbf{0.604} & \textbf{0.660} \\
\bottomrule
\multicolumn{5}{l}{\footnotesize PL~: PiecewiseLinear embeddings.}
\end{tabular}
\end{table}

L'ordre de performance est conforme aux prédictions de l'article~:
TabM-packed~$<$~TabM-mini~$<$~TabM~$<$~TabM-mini$^\dagger$~$<$~TabM$^\dagger$.
Le gain des embeddings PiecewiseLinear est le plus significatif ($+0.011$~QWK),
confirmant que la modélisation explicite des non-linéarités numériques est le
facteur clé sur ce jeu de données. L'optimisation des offsets apporte
systématiquement $+0.05$~QWK sans aucun ré-entraînement.

\subsubsection{Sweep bayésien}

Un sweep bayésien (Weights \& Biases, 30~runs) est conduit sur la variante
TabM$^\dagger$ pour optimiser les hyperparamètres
$\eta$, $\lambda$, \textit{dropout}, $B$, $d_e$, $n_{\text{blocs}}$.
Le meilleur run (\texttt{wise-sweep-22}) atteint un QWK de~$\mathbf{0.6608}$
avec la configuration~: $\eta=1.12\times10^{-3}$, $\lambda=3.37\times10^{-4}$,
\textit{dropout}~$=0.229$, $B=48$, $d_e=32$, $n_{\text{blocs}}=3$.
La substitution de $d_e=16$ (valeur par défaut) par $d_e=32$ est le seul
changement substantiel, expliquant l'essentiel du gain de~$+0.001$.

\begin{table}[htbp]
\caption{Top-5 runs du sweep bayésien (TabM$^\dagger$)}
\label{tab:sweep}
\centering
\begin{tabular}{lcccc}
\toprule
\textbf{Run} & \textbf{QWK\textsubscript{+off.}} & $\boldsymbol{\eta}$ &
$\boldsymbol{d_e}$ & $n_b$ \\
\midrule
wise-sweep-22    & \textbf{0.6608} & $1.12\times10^{-3}$ & 32 & 3 \\
mild-sweep-18    & 0.6605 & $5.67\times10^{-4}$ & 32 & 3 \\
blooming-sweep-10 & 0.6591 & $8.23\times10^{-4}$ & 8  & 3 \\
avid-sweep-7     & 0.6590 & $1.46\times10^{-3}$ & 32 & 2 \\
lyric-sweep-24   & 0.6590 & $7.48\times10^{-4}$ & 8  & 3 \\
\bottomrule
\end{tabular}
\end{table}

\subsubsection{Ensemble stacking}

Afin d'explorer si la combinaison de modèles hétérogènes peut dépasser
TabM$^\dagger$ seul, nous entraînons un ensemble par \textit{stacking}
à deux niveaux en validation croisée stratifiée ($k=5$)~:
\begin{itemize}
  \item \textbf{Niveau~0}~: TabM$^\dagger$ (hyperparamètres sweep optimaux),
        XGBoost ($n_\text{est}=800$, $\eta=0.05$, \textit{depth}=6),
        LightGBM ($n_\text{est}=800$, $\eta=0.05$, 63~feuilles)~;
  \item \textbf{Niveau~1}~: moyenne pondérée optimisée ($w_{\text{TabM}}=0.4$,
        $w_{\text{XGB}}=0.3$, $w_{\text{LGB}}=0.3$) suivie d'une optimisation
        des offsets.
\end{itemize}

\begin{table}[htbp]
\caption{Résultats du stacking en validation croisée OOF ($k=5$)}
\label{tab:stacking}
\centering
\begin{tabular}{lcc}
\toprule
\textbf{Modèle} & \textbf{QWK\textsubscript{brut}} &
\textbf{QWK\textsubscript{+off.}} \\
\midrule
TabM$^\dagger$ (OOF)   & 0.594 & 0.624 \\
XGBoost (OOF)          & 0.611 & 0.635 \\
LightGBM (OOF)         & 0.607 & 0.637 \\
\textbf{Stacking}      & \textbf{0.606} & \textbf{0.637} \\
\midrule
TabM$^\dagger$ (ablation, full train) & 0.604 & \textbf{0.660} \\
\bottomrule
\end{tabular}
\end{table}

Le stacking ne dépasse pas TabM$^\dagger$ entraîné sur la totalité des
données d'entraînement ($0.637$ vs $0.660$). Deux facteurs l'expliquent~:
(1)~TabM$^\dagger$ perd~$\approx 0.036$~QWK en mode OOF (entraîné sur
$80\,\%$ des données seulement, avec \textit{early stopping})~;
(2)~XGBoost et LightGBM sont très fortement corrélés entre eux
($r=0.986$), réduisant la diversité effective de l'ensemble. La corrélation
entre TabM$^\dagger$ et les modèles boosting est quasi nulle ($r\approx0.01$),
ce qui confirme leur complémentarité théorique, mais TabM est trop affaibli
par le régime OOF pour en tirer parti.

\subsection{Ce qui n'a pas fonctionné}

\textbf{[DOCUMENTER LES ÉCHECS]} : Quelles approches ou configurations n'ont pas donné les
résultats escomptés? Quelles hypothèses ont été infirmées?

Concernant la partie TabM, le \textbf{stacking} n'a pas permis de dépasser TabM$^\dagger$
seul~: l'affaiblissement de TabM en régime OOF ($-0.036$~QWK) annule le bénéfice attendu
de la diversité. Des runs supplémentaires du sweep bayésien au-delà de 30 n'apportaient
plus de gain mesurable, les réglages par défaut étant déjà quasi-optimaux.

\subsection{Possibilités d'amélioration}

\textbf{[SUGGESTIONS]} : Quelles pistes pourraient être explorées avec plus de temps ou de
ressources?

Pour la partie TabM, plusieurs pistes seraient envisageables~: (1)~augmenter le nombre
d'epochs par fold lors du stacking (actuellement limité par l'\textit{early stopping} sur
80\,\% des données)~; (2)~introduire un modèle tabm-packed à architecture différente pour
augmenter la diversité de l'ensemble~; (3)~explorer des valeurs de $k>32$ ou des architectures
à plus de blocs pour les variantes non-embedding~; (4)~combiner TabM$^\dagger$ avec des
descripteurs issus de TabDPT pour un stacking inter-paradigmes.

% ---------------------------------------------------------------
% 6. Conclusion
% ---------------------------------------------------------------
\section{Conclusion}

Ce projet a présenté une exploration comparative d'architectures d'apprentissage profond
modernes appliquées à l'évaluation du risque ordinal en assurance-vie. En utilisant le
jeu de données \textit{Prudential Life Insurance Assessment}, nous avons confronté des
modèles de \textit{gradient boosting} bien établis (XGBoost, LightGBM, CatBoost) à des
architectures récentes conçues pour les données tabulaires (TabPFN-2.5, TabM, TabDPT).

Nos expériences montrent que les architectures récentes dédiées aux données tabulaires
surpassent les MLP standards tout en restant compétitives face aux méthodes de
\textit{gradient boosting}. En particulier,
\textbf{TabM$^\dagger$ avec embeddings PiecewiseLinear atteint un QWK de 0,6608 (meilleur
sweep W\&B) contre 0,6320 pour le MLP de référence, soit un gain de $+0.029$}.
L'étude d'ablation révèle que \textbf{les embeddings PiecewiseLinear apportent le gain le
plus important ($+0.011$~QWK), suivis par la régularisation BatchEnsemble ($+0.008$~QWK),
tandis que l'architecture \textit{packed} n'apporte pas de bénéfice sur ce jeu de données}.
L'approche de \textbf{stacking TabM + XGBoost + LightGBM (QWK~=~0.6369) n'a pas dépassé
TabM seul, la forte corrélation entre les prédicteurs (r~=~0.986) annulant tout effet de
diversité}.

Ces résultats suggèrent que \textbf{les embeddings numériques spécialisés (PiecewiseLinear)
et la régularisation par BatchEnsemble sont les facteurs discriminants qui permettent aux
architectures de type TabM de concurrencer, voire dépasser, les méthodes de \textit{gradient
boosting} sur des données tabulaires hétérogènes à variable ordinale}, ce qui contribue à la compréhension
du potentiel et des limites de l'apprentissage profond sur les données tabulaires dans un
contexte applicatif réel.

% ---------------------------------------------------------------
% 7. Utilisation de l'IA générative
% ---------------------------------------------------------------
\section{Utilisation de l'IA générative}

\textbf{[COMPLÉTER CETTE SECTION CONFORMÉMENT AUX DIRECTIVES]}

Dans le cadre de ce projet, nous avons eu recours à des outils d'IA générative pour les
usages suivants~:

\begin{itemize}
  \item \textbf{Génération de code}~: \textbf{[DÉCRIRE]}, notamment pour les fonctions
        \texttt{[NOMS DES FONCTIONS]}. Exemples de prompts utilisés~: \textit{``[PROMPT 1]''},
        \textit{``[PROMPT 2]''}.
  \item \textbf{Lissage grammatical}~: Les sections \textbf{[NOMS DES SECTIONS]} ont été
        relues et corrigées avec l'aide de \textbf{[OUTIL]}.
\end{itemize}

% OU, si aucun usage :
% Nous n'avons pas eu recours à des outils d'IA générative dans la rédaction de ce rapport
% ni dans la génération du code associé à ce projet.

% ---------------------------------------------------------------
% Références
% ---------------------------------------------------------------
\bibliographystyle{IEEEtran}
\bibliography{references}

\end{document}
